\ifdefined\gkver\else\def\gkver{student}\fi
\documentclass[\gkver]{gaokaozhenti}
\begin{document}

\chapter{2010年福建理综}
%% paperType: 理综物理部分

\begin{enumerate}
\item
%% number: 13
%% paperName: 2010年福建理综第 13 题 6 分
%% typeId: 1
%% type: 单选题
%% chapter: 电磁感应
%% point: 远距离输电
%% method: 无
%% score: 6
%% degree: 650
%% duplicateId: 0
%% body:
中国已投产运行的 $1000\UkV$ 特高压输电是目前世界上电压最高的输电工程。假设甲、乙两地原来用 $500\UkV$ 的超高压输电，输电线上损耗的电功率为 $P$。保持输送电功率和输电线电阻都不变的条件下，现改用 $1000\UkV$ 特高压输电，若不考虑其他因素的影响，则输电线上损耗的电功率将变为\xzanswer{A}
\fourchoices[answer=A]
{$\frac{P}{4}$}
{$\frac{P}{2}$}
{$2P$}
{$4P$}
%% answer: A
%% memo:
\memoanswer{
由 $P=UI$ 可知输出电压由 $500\UkV$ 升高到 $1000\UkV$ 时，电路中的电流将减为原来的一半；由 $P_{\text{损}}=I^{2}R$ 可知电路中损耗的功率变为原来的 $\frac{1}{4}$，选项 A 正确。
}

\item
%% number: 14
%% paperName: 2010年福建理综第 14 题 6 分
%% typeId: 1
%% type: 单选题
%% chapter: 曲线运动
%% point: 天体运动
%% method: 无
%% score: 6
%% degree: 450
%% duplicateId: 0
%% body:
火星探测项目是我国继神舟载人航天工程、嫦娥探月工程之后又一个重大太空探索项目。假设火星探测器在火星表面附近圆轨道运行的周期为 $T_{1}$，神舟飞船在地球表面附近的圆形轨道运行周期为 $T_{2}$，火星质量与地球质量之比为 $p$，火星半径与地球半径之比为 $q$，则 $T_{1}$ 与 $T_{2}$ 之比为\xzanswer{D}
\fourchoices[answer=D]
{$\sqrt{pq^{3}}$}
{$\sqrt{\frac{1}{pq^{3}}}$}
{$\sqrt{\frac{p}{q^{3}}}$}
{$\sqrt{\frac{q^{3}}{p}}$}
%% answer: D
%% memo:
\memoanswer{
设中心天体的质量为 $M$、半径为 $R$，当航天器在星球表面飞行时，由 $G\frac{Mm}{R^{2}}=m\left(\frac{2\pi}{T}\right)^{2}R$ 和 $M=\rho V=\rho\cdot\frac{4}{3}\pi R^{3}$ 解得 $\rho=\frac{3\pi}{GT^{2}}$，即 $T=\sqrt{\frac{3\pi}{\rho G}}\propto\sqrt{\frac{1}{\rho}}$；又因为 $\rho=\frac{M}{V}=\frac{M}{\frac{4}{3}\pi R^{3}}\propto\frac{M}{R^{3}}$，所以 $T\propto\sqrt{\frac{R^{3}}{M}}$，$\frac{T_{1}}{T_{2}}=\sqrt{\frac{q^{3}}{p}}$。选项 D 正确。
}

\item
%% number: 15
%% paperName: 2010年福建理综第 15 题 6 分
%% typeId: 1
%% type: 单选题
%% chapter: 机械波
%% point: 波的多解性
%% method: 无
%% score: 6
%% degree: 450
%% duplicateId: 0
%% body:
一列简谐横波在 $t=0$ 时刻的波形如图中的实线所示，$t=0.02\Us$ 时刻的波形如图中虚线所示。若该波的周期 $T$ 大于 $0.02\Us$，则该波的传播速度可能是\xzanswer{B}

\onepicture[align=right, w=6.5cm]{figs/15.png}

\fourchoices[answer=B]
{$2\Ums$}
{$3\Ums$}
{$4\Ums$}
{$5\Ums$}
%% answer: B
%% memo:
\memoanswer{
解法一：质点振动法

（1）设波向右传播，则在 0 时刻 $x=4\Um$ 处的质点往上振动，设经历 $\Delta t$ 时间时质点运动到波峰的位置，则 $\Delta t=\left(\frac{1}{4}+n\right)T$，即 $T=\frac{4\Delta t}{4n+1}=\frac{0.08}{4n+1}$。当 $n=0$ 时，$T=0.08\Us>0.02\Us$，符合要求，此时 $v=\frac{\lambda}{T}=\frac{0.08}{0.08}=1\Ums$；当 $n=1$ 时，$T=0.016\Us<0.02\Us$，不符合要求。

（2）设波向左传播，则在 0 时刻 $x=4\Um$ 处的质点往下振动，设经历 $\Delta t$ 时间时质点运动到波峰的位置，则 $\Delta t=\left(\frac{3}{4}+n\right)T$，即 $T=\frac{4\Delta t}{4n+3}=\frac{0.08}{4n+3}$。当 $n=0$ 时，$T=\frac{0.08}{3}\Us>0.02\Us$，符合要求，此时 $v=\frac{\lambda}{T}=\frac{0.08}{\frac{0.08}{3}}=3\Ums$；当 $n=1$ 时，$T=\frac{0.08}{7}\Us<0.02\Us$，不符合要求。

综上所述，只有选项 B 正确。

解法二：波的平移法

（1）设波向右传播，只有当波传播的距离为 $\Delta x=0.02+0.08n$ 时，实线才会和虚线重合，即 0 时刻的波形才会演变成 $0.02\Us$ 时的波形，$\Delta t=0.02\Us$，所以 $v=\frac{\Delta x}{\Delta t}=\frac{0.02+0.08n}{0.02}=1+4n$，$T=\frac{\lambda}{v}=\frac{0.08}{4n+1}$。当 $n=0$ 时，$T=0.08\Us>0.02\Us$，符合要求，此时 $v=1+4n=1\Ums$；当 $n=1$ 时，$T=0.016\Us<0.02\Us$，不符合要求。

（2）设波向左传播，只有当波传播的距离为 $\Delta x=0.06+0.08n$ 时，实线才会和虚线重合，即 0 时刻的波形才会演变成 $0.02\Us$ 时的波形，$\Delta t=0.02\Us$，所以 $v=\frac{\Delta x}{\Delta t}=\frac{0.06+0.08n}{0.02}=3+4n$，$T=\frac{\lambda}{v}=\frac{0.08}{4n+3}$。当 $n=0$ 时，$T=\frac{0.08}{3}\Us>0.02\Us$，符合要求，此时 $v=3+4n=3\Ums$；当 $n=1$ 时，$T=\frac{0.08}{7}\Us<0.02\Us$，不符合要求。

综上所述，只有选项 B 正确。
}

\item
%% number: 16
%% paperName: 2010年福建理综第 16 题 6 分
%% typeId: 1
%% type: 单选题
%% chapter: 运动和力
%% point: 牛顿定律中的图象
%% method: 无
%% score: 6
%% degree: 450
%% duplicateId: 0
%% body:
质量为 $2\Ukg$ 的物体静止在足够大的水平地面上，物体与地面间的动摩擦因数为 0.2，最大静摩擦力与滑动摩擦力大小视为相等。从 $t=0$ 时刻开始，物体受到方向不变、大小呈周期性变化的水平拉力 $F$ 的作用，$F$ 随时间 $t$ 的变化规律如图所示。重力加速度 $g$ 取 $10\Umsq$，则物体在 $t=0$ 到 $t=12\Us$ 这段时间的位移大小为\xzanswer{B}

\onepicture[align=right, w=5.5cm]{figs/16.png}

\fourchoices[answer=B]
{$18\Um$}
{$54\Um$}
{$72\Um$}
{$198\Um$}
%% answer: B
%% memo:
\memoanswer{
拉力只有大于最大静摩擦力时，物体才会由静止开始运动。

0-3s 时：$F=f_{\max}$，物体保持静止，$s_{1}=0$；

3-6s 时：$F>f_{\max}$，物体由静止开始做匀加速直线运动，$a=\frac{F-f}{m}=\frac{8-4}{2}\Umsq=2\Umsq$，$v=at=6\Ums$，$s_{2}=\frac{1}{2}at^{2}=\frac{1}{2}\times2\times3^{2}\Um=9\Um$；

6-9s 时：$F=f$，物体做匀速直线运动，$s_{3}=vt=6\times3=18\Um$；

9-12s 时：$F>f$，物体以 $6\Ums$ 为初速度，以 $2\Umsq$ 为加速度继续做匀加速直线运动，$s_{4}=vt+\frac{1}{2}at^{2}=6\times3+\frac{1}{2}\times2\times3^{2}=27\Um$；

所以 0-12s 内物体的位移为 $s=s_{1}+s_{2}+s_{3}+s_{4}=54\Um$，选项 B 正确。
}

\item
%% number: 17
%% paperName: 2010年福建理综第 17 题 6 分
%% typeId: 1
%% type: 单选题
%% chapter: 运动和力
%% point: 变加速直线运动
%% method: 无
%% score: 6
%% degree: 450
%% duplicateId: 0
%% body:
如图 \ref{2010福建理综17a} 所示，质量不计的弹簧竖直固定在水平面上，$t=0$ 时刻，将一金属小球从弹簧正上方某一高度处由静止释放，小球落到弹簧上压缩弹簧到最低点，然后又被弹起离开弹簧，上升到一定高度后再下落，如此反复。通过安装在弹簧下端的压力传感器，测出这一过程弹簧弹力 $F$ 随时间 $t$ 变化的图像如图 \ref{2010福建理综17b} 如示，则\xzanswer{C}

\twopicture[numA=甲, numB=乙, labelA=2010福建理综17a, labelB=2010福建理综17b, wA=4.2cm, wB=6.5cm]
{figs/17a.png}{figs/17b.png}

\fourchoices[answer=C]
{$t_{1}$ 时刻小球动能最大}
{$t_{2}$ 时刻小球动能最大}
{$t_{2}-t_{3}$ 这段时间内，小球的动能先增加后减少}
{$t_{2}-t_{3}$ 这段时间内，小球增加的动能等于弹簧减少的弹性势能}
%% answer: C
%% memo:
\memoanswer{
小球在接触弹簧之前做自由落体。碰到弹簧后先做加速度不断减小的加速运动，当加速度为 0 即重力等于弹簧弹力时速度达到最大值，而后往下做加速度不断增大的减速运动，与弹簧接触的整个下降过程，小球的动能和重力势能转化为弹簧的弹性势能。上升过程恰好与下降过程互逆。由乙图可知 $t_{1}$ 时刻开始接触弹簧；$t_{2}$ 时刻弹力最大，小球处在最低点，动能最小；$t_{3}$ 时刻小球往上运动恰好要离开弹簧；$t_{2}-t_{3}$ 这段时间内，小球先加速后减速，动能先增加后减少，弹簧的弹性势能转化为小球的动能和重力势能。
}

\item
%% number: 18
%% paperName: 2010年福建理综第 18 题 6 分
%% typeId: 1
%% type: 单选题
%% chapter: 电场
%% point: 电场叠加
%% method: 微元
%% score: 6
%% degree: 250
%% duplicateId: 0
%% body:
物理学中有些问题的结论不一定必须通过计算才能验证，有时只需通过一定的分析就可以判断结论是否正确。如图所示为两个彼此平行且共轴的半径分别为 $R_{1}$ 和 $R_{2}$ 的圆环，两圆环上的电荷量均为 $q$（$q>0$），而且电荷均匀分布。两圆环的圆心 O$_{1}$ 和 O$_{2}$ 相距为 $2a$，连线的中点为 O，轴线上的 A 点在 O 点右侧与 O 点相距为 $r$（$r<a$）。试分析判断下列关于 A 点处电场强度大小 $E$ 的表达式（式中 $k$ 为静电力常量）正确的是\xzanswer{D}

\onepicture[align=right, w=7cm]{figs/18.png}

\fourchoices[answer=D]
{$E=\left|\frac{kqR_{1}}{R_{1}^{2}+(a+r)^{2}}-\frac{kqR_{2}}{R_{2}^{2}+(a-r)^{2}}\right|$}
{$E=\left|\frac{kqR_{1}}{[R_{1}^{2}+(a+r)^{2}]^{3/2}}-\frac{kqR_{2}}{[R_{2}^{2}+(a-r)^{2}]^{3/2}}\right|$}
{$E=\left|\frac{kq(a+r)}{R_{1}^{2}+(a+r)^{2}}-\frac{kq(a-r)}{R_{2}^{2}+(a-r)^{2}}\right|$}
{$E=\left|\frac{kq(a+r)}{[R_{1}^{2}+(a+r)^{2}]^{3/2}}-\frac{kq(a-r)}{[R_{2}^{2}+(a-r)^{2}]^{3/2}}\right|$}
%% answer: D
%% memo:
\memoanswer{
当 $r=0$ 时，点位于圆心 O 处，可以把 O$_{1}$、O$_{2}$ 两个带电圆环均等效成两个位于圆心处的点电荷，根据场强的叠加容易知道，此时总场强 $E=0$，将 $r=0$ 代入各选项，排除 A、B 选项；当 $r=a$ 时，A 点位于圆心 O$_{2}$ 处，带电圆环 O$_{2}$ 由于对称性在 A 点的电场为 0，根据微元法可以求得此时的总场强为 $E=E_{1}=\frac{2kqa}{[R_{1}^{2}+4a^{2}]^{3/2}}$，将 $r=a$ 代入 C、D 选项可以排除 C。
}

\item
%% number: 19
%% paperName: 2010年福建理综第 19 题 6 分
%% typeId: 4
%% type: 实验题
%% chapter: 电路
%% point: 电路实验
%% method: 无
%% score: 6
%% degree: 450
%% duplicateId: 0
%% body:
（18 分）\textbf{（1）}（6 分）某同学利用“插针法”测定玻璃的折射率，所用的玻璃砖两面平行。正确操作后，作出的光路图及测出的相关角度如图所示。

①此玻璃的折射率计算式为 $n=$\tkanswer[2.5cm]{$\frac{\cos\theta_{1}}{\cos\theta_{2}}$}（用图中的 $\theta_{1}$、$\theta_{2}$ 表示）；

②如果有几块宽度大小不同的平行玻璃砖可供选择，为了减小误差，应选用宽度\tkanswer[1cm]{大}（填“大”或“小”）的玻璃砖来测量。

\onepicture[align=right, w=4.5cm]{figs/19a.png}

\textbf{（2）}（6 分）某实验小组研究橡皮筋伸长与所受拉力的关系。实验时，将原长约 $200\Umm$ 的橡皮筋上端固定，在竖直悬挂的橡皮筋下端逐一增挂钩码（质量均为 $20\Ug$），每增挂一只钩码均记下对应的橡皮筋伸长量；当挂上 10 只钩码后，再逐一把钩码取下，每取下一只钩码，也记下对应的橡皮筋伸长量。根据测量数据，作出增挂钩码和减挂钩码时的橡皮筋伸长量 $\Delta l$ 与拉力 $F$ 关系的图像如图所示。从图像中可以得出\tkanswer[1cm]{D}。（填选项前的字母）

（A）增挂钩码时 $\Delta l$ 与 $F$ 成正比，而减挂钩码时 $\Delta l$ 与 $F$ 不成正比

（B）当所挂钩码数相同时，增挂钩码时橡皮筋的伸长量比减挂钩码时的大

（C）当所挂钩码数相同时，增挂钩码时橡皮筋的伸长量与减挂钩码时的相等

（D）增挂钩码时所挂钩码数过多，导致橡皮筋超出弹性限度

\onepicture[align=right, w=5cm]{figs/19b.png}

\textbf{（3）}（6 分）如图所示是一些准备用来测量待测电阻 $R_{x}$ 阻值的实验器材，器材及其规格列表如下：

\begin{center}
\begin{tblr}{|c|c|}
\hline
器材 & 规格 \\
\hline
待测电阻 $R_{x}$ & 阻值在 $900\UO\sim1000\UO$ 之间 \\
电源 $E$ & 具有一定内阻，电动势约 $9.0\UV$ \\
电压表 V$_{1}$ & 量程 $2.0\UV$，内阻 $r_{1}=1000\UO$ \\
电压表 V$_{2}$ & 量程 $5.0\UV$，内阻 $r_{2}=2500\UO$ \\
电流表 A & 量程 $3.0\UA$，内阻 $r=0.10\UO$ \\
滑动变阻器 $R$ & 最大阻值约 $100\UO$，额定电流 $0.5\UA$ \\
开关 S、导线若干 & \\
\hline
\end{tblr}
\end{center}

为了能正常进行测量并尽可能减小测量误差，实验要求测量时电表的读数大于其量程的一半，而且调节滑动变阻器能使电表读数有较明显的变化。请用实线代替导线，在所给的实验器材图中选择若干合适的器材，连成满足要求的测量 $R_{x}$ 阻值的电路。

\onepicture[align=right, w=6cm]{figs/19c.png}

%% answer: 
\jdanswer{
\begin{enumerate}
	\item
	① $\frac{\cos\theta_{1}}{\cos\theta_{2}}$；② 大。

	\item
	D。

	\item
	连线如图所示。

\end{enumerate}
}
%% memo:
\memoanswer{
（1）由折射定律 $n=\frac{\sin i}{\sin r}=\frac{\sin(90^{\circ}-\theta_{1})}{\sin(90^{\circ}-\theta_{2})}=\frac{\cos\theta_{1}}{\cos\theta_{2}}$。应选宽一些的玻璃砖。

（2）本小题解题的关键在于弄清为什么两个图象会错开。实验中由于增挂钩码时所挂钩码数过多，导致橡皮筋超出弹性限度，弹簧无法及时回复原长，两条图象才会出现差异。

（3）要准确就需要对实验原理（伏安法、变阻器接法）、误差分析有深刻理解。由于待测电阻约 $1000\UO$，滑动变阻器最大电阻为 $100\UO$，要使电表读数发生变化，滑动变阻器必须使用分压式接法。当电源电动势全部加在待测电阻上时，流过的电流约为 $I=\frac{E}{R_{x}}=\frac{9}{1000}=9\UmA\ll3.0\UA$，显然不能用电流表来测电流，而应该把其中一个电压表 V$_{1}$ 当电流表来使用。根据比值法容易判断出待测电阻是大电阻，伏安法应使用内接法。

\onepicture[w=6cm]{figs/19_an.png}
}

\item
%% number: 20
%% paperName: 2010年福建理综第 20 题 15 分
%% typeId: 6
%% type: 计算题
%% chapter: 磁场
%% point: 带电粒子在复合场中的运动
%% method: 无
%% score: 15
%% degree: 650
%% duplicateId: 0
%% body:
如图所示的装置，左半部为速度选择器，右半部为匀强的偏转电场。一束同位素离子流从狭缝 S$_{1}$ 射入速度选择器，能够沿直线通过速度选择器并从狭缝 S$_{2}$ 射出的离子，又沿着与电场垂直的方向，立即进入场强大小为 $E$ 的偏转电场，最后打在照相底片 D 上。已知同位素离子的电荷量为 $q$（$q>0$），速度选择器内部存在着相互垂直的场强大小为 $E_{0}$ 的匀强电场和磁感应强度大小为 $B_{0}$ 的匀强磁场，照相底片 D 与狭缝 S$_{1}$、S$_{2}$ 的连线平行且距离为 $L$，忽略重力的影响。
\begin{enumerate}
	\item
	求从狭缝 S$_{2}$ 射出的离子速度 $v_{0}$ 的大小；

	\item
	若打在照相底片上的离子在偏转电场中沿速度 $v_{0}$ 方向飞行的距离为 $x$，求出 $x$ 与离子质量 $m$ 之间的关系式（用 $E_{0}$、$B_{0}$、$E$、$q$、$m$、$L$ 表示）。

\end{enumerate}

\onepicture[align=right, w=8cm]{figs/20.png}

%% answer: 
\jdanswer{
\begin{enumerate}
	\item
	$\frac{E_{0}}{B_{0}}$

	\item
	$x=\frac{E_{0}}{B_{0}}\sqrt{\frac{2mL}{qE}}$

\end{enumerate}
}
%% memo:
\memoanswer{
（1）能从速度选择器射出的离子满足 $qE_{0}=qv_{0}B_{0}$，解得 $v_{0}=\frac{E_{0}}{B_{0}}$。

（2）离子进入匀强偏转电场 $E$ 后做类平抛运动，则 $x=v_{0}t$，$L=\frac{1}{2}at^{2}$，由牛顿第二定律 $qE=ma$，解得 $x=\frac{E_{0}}{B_{0}}\sqrt{\frac{2mL}{qE}}$。
}

\item
%% number: 21
%% paperName: 2010年福建理综第 21 题 19 分
%% typeId: 6
%% type: 计算题
%% chapter: 电磁感应
%% point: 双杆问题
%% method: 无
%% score: 19
%% degree: 450
%% duplicateId: 0
%% body:
如图所示，两条平行的光滑金属导轨固定在倾角为 $\theta$ 的绝缘斜面上，导轨上端连接一个定值电阻。导体棒 a 和 b 放在导轨上，与导轨垂直并良好接触。斜面上水平虚线 PQ 以下区域内，存在着垂直穿过斜面向上的匀强磁场。现对 a 棒施以平行导轨斜向上的拉力，使它沿导轨匀速向上运动，此时放在导轨下端的 b 棒恰好静止。当 a 棒运动到磁场的上边界 PQ 处时，撤去拉力，a 棒将继续沿导轨向上运动一小段距离后再向下滑动，此时 b 棒已滑离导轨。当 a 棒再次滑回到磁场边界 PQ 处时，又恰能沿导轨匀速向下运动。已知 a 棒、b 棒和定值电阻的阻值均为 $R$，b 棒的质量为 $m$，重力加速度为 $g$，导轨电阻不计。求：
\begin{enumerate}
	\item
	a 棒在磁场中沿导轨向上运动的过程中，a 棒中的电流强度 $I_{a}$ 与定值电阻 $R$ 中的电流强度 $I_{R}$ 之比；

	\item
	a 棒质量 $m_{a}$；

	\item
	a 棒在磁场中沿导轨向上运动时所受的拉力 $F$。

\end{enumerate}

\onepicture[align=right, w=5cm]{figs/21.png}

%% answer: 
\jdanswer{
\begin{enumerate}
	\item
	$2:1$

	\item
	$m_{a}=\frac{3}{2}m$

	\item
	$\frac{7}{2}mg\sin\theta$

\end{enumerate}
}
%% memo:
\memoanswer{
（1）a 棒沿导轨向上运动时，a 棒、b 棒及电阻 $R$ 中的电流分别为 $I_{a}$、$I_{b}$、$I_{R}$，有 $I_{R}R=I_{b}R_{1}$，$I_{a}=I_{b}+I_{R}$，解得 $\frac{I_{a}}{I_{R}}=\frac{2}{1}$。

（2）由于 a 棒在 PQ 上方滑动过程中机械能守恒，因而在磁场中向上滑动的速度大小 $v_{1}$ 与在磁场中向下滑动的速度大小 $v_{2}$ 相等，即 $v_{1}=v_{2}=v$。设磁场的磁感应强度为 $B$，导体棒长为 $L$，a 棒在磁场中运动时产生的感应电动势为 $E=BLv$。当 a 棒向上滑动时 $I_{b}=\frac{E}{3R}$，$I_{b}BL=mg\sin\theta$；当 a 棒向下滑动时 $I_{a}'=\frac{E}{2R}$，$I_{a}'BL=m_{a}g\sin\theta$，解得 $m_{a}=\frac{3}{2}m$。

（3）由题知导体棒 a 沿斜面向上运动时，所受拉力 $F=I_{a}BL+m_{a}g\sin\theta=\frac{7}{2}mg\sin\theta$。
}

\item
%% number: 22
%% paperName: 2010年福建理综第 22 题 20 分
%% typeId: 6
%% type: 计算题
%% chapter: 运动和力
%% point: 传送带
%% method: 无
%% score: 20
%% degree: 450
%% duplicateId: 0
%% body:
如图所示，物体 A 放在足够长的木板 B 上，木板 B 静止于水平面。$t=0$ 时，电动机通过水平细绳以恒力 $F$ 拉木板 B，使它做初速度为零、加速度 $a_{B}=1.0\Umsq$ 的匀加速直线运动。已知 A 的质量 $m_{A}$ 和 B 的质量 $m$ 均为 $2.0\Ukg$，A、B 之间的动摩擦因数 $\mu_{1}=0.05$，B 与水平面之间的动摩擦因数 $\mu_{2}=0.1$，最大静摩擦力与滑动摩擦力大小视为相等，重力加速度 $g$ 取 $10\Umsq$。求：
\begin{enumerate}
	\item
	物体 A 刚运动时的加速度 $a_{A}$；

	\item
	$t=1.0\Us$ 时，电动机的输出功率 $P$；

	\item
	若 $t=1.0\Us$ 时，将电动机的输出功率立即调整为 $P'=5\UW$，并在以后的运动过程中始终保持这一功率不变，$t=3.8\Us$ 时物体 A 的速度为 $1.2\Ums$。则在 $t=1.0\Us$ 到 $t=3.8\Us$ 这段时间内木板 B 的位移为多少？

\end{enumerate}

\onepicture[align=right, w=10cm]{figs/22.png}

%% answer: 
\jdanswer{
\begin{enumerate}
	\item
	$0.5\Umsq$

	\item
	$7\UW$

	\item
	$3.03\Um$

\end{enumerate}
}
%% memo:
\memoanswer{
（1）物体 A 在水平方向上受到向右的摩擦力，由牛顿第二定律得 $\mu_{1}m_{A}g=m_{A}a_{A}$，代入数据解得 $a_{A}=0.5\Umsq$。

（2）$t=1.0\Us$ 时，木板 B 的速度大小为 $v=a_{B}t=1\Ums$。木板 B 所受拉力 $F$，由牛顿第二定律有 $F-\mu_{1}m_{A}g-\mu_{2}(m_{A}+m_{B})g=m_{B}a_{B}$，解得 $F=7\UN$，电动机输出功率 $P=Fv=7\UW$。

（3）电动机的输出功率调整为 $5\UW$ 时，设细绳对木板 B 的拉力为 $F'$，则 $P'=F'v$，解得 $F'=5\UN$，木板 B 受力满足 $F'-\mu_{1}m_{A}g-\mu_{2}(m_{A}+m_{B})g=0$，所以木板 B 将做匀速直线运动，而物体 A 则继续在 B 上做匀加速直线运动，直到 A、B 速度相等。设这一过程时间为 $t'$，有 $v=a_{A}(t_{1}+t')$，这段时间内木板 B 的位移 $s_{1}=vt'$。A、B 速度相等后，由于 $F'>\mu_{2}(m_{A}+m_{B})g$ 且电动机输出功率恒定，A、B 将一起做加速度逐渐减小的变加速运动，由动能定理得 $P'(t_{2}-t'-t_{1})-\mu_{2}(m_{A}+m_{B})gs_{2}=\frac{1}{2}(m_{A}+m_{B})v_{A}^{2}-\frac{1}{2}(m_{A}+m_{B})v_{1}^{2}$，由以上各式代入数据解得木板 B 在 $t=1.0\Us$ 到 $3.8\Us$ 这段时间内的位移 $s=s_{1}+s_{2}=3.03\Um$。
}

\item
%% number: 28
%% paperName: 2010年福建理综第 28 题 6 分
%% typeId: 1
%% type: 单选题
%% chapter: 气体性质
%% point: 热力学第一定律
%% method: 无
%% score: 6
%% degree: 650
%% duplicateId: 0
%% body:
【物理——选修 3-3】（本题共有两小题，每小题 6 分，共 12 分。每小题只有一个选项符合题意）

\textbf{（1）}1859 年麦克斯韦从理论上推导出了气体分子速率的分布规律，后来有许多实验验证了这一规律。若以横坐标 $v$ 表示分子速率，纵坐标 $f(v)$ 表示各速率区间的分子数占总分子数的百分比。下面各幅图中能正确表示某一温度下气体分子速率分布规律的是\tkanswer[1cm]{D}。（填选项前的字母）

\fourchoices[ispicture=true, h=1.5cm]
{figs/28a.png}
{figs/28b.png}
{figs/28c.png}
{figs/28d.png}

\textbf{（2）}如图所示，一定质量的理想气体密封在绝热（即与外界不发生热交换）容器中，容器内装有一可以活动的绝热活塞。今对活塞施以一竖直向下的压力 $F$，使活塞缓慢向下移动一段距离后，气体的体积减小。若忽略活塞与容器壁间的摩擦力，则被密封的气体\xzanswer{C}。（填选项前的字母）

\onepicture[align=right, w=2.2cm]{figs/28e.png}

\fourchoices[answer=C]
{温度升高，压强增大，内能减少}
{温度降低，压强增大，内能减少}
{温度升高，压强增大，内能增加}
{温度降低，压强减小，内能增加}
%% answer: C
%% memo:
\memoanswer{
（1）图象突出中间多、两边少的特点，答案选 D。

（2）外力 $F$ 做正功，$W>0$；绝热，$Q=0$；由热力学第一定律 $\Delta U=Q+W>0$，内能增加，温度升高；另外，由 $\frac{pV}{T}=C$ 可以判断出压强增大。故选 C。
}

\item
%% number: 29
%% paperName: 2010年福建理综第 29 题 6 分
%% typeId: 1
%% type: 单选题
%% chapter: 运动和力
%% point: 动量守恒定律
%% method: 无
%% score: 6
%% degree: 650
%% duplicateId: 0
%% body:
【物理——选修 3-5】（本题共有两小题，每小题 6 分，共 12 分。每小题只有一个选项符合题意）

\textbf{（1）}$\ce{^{14}C}$ 测年法是利用 $\ce{^{14}C}$ 衰变规律对古生物进行年代测定的方法。若以横坐标 $t$ 表示时间，纵坐标 $m$ 表示任意时刻 $\ce{^{14}C}$ 的质量，$m_{0}$ 为 $t=0$ 时 $\ce{^{14}C}$ 的质量。下面四幅图中能正确反映 $\ce{^{14}C}$ 衰变规律的是\tkanswer[1cm]{C}。（填选项前的字母）

\fourchoices[ispicture=true, h=1.5cm]
{figs/29a.png}
{figs/29b.png}
{figs/29c.png}
{figs/29d.png}

\textbf{（2）}如图所示，一个木箱原来静止在光滑水平面上，木箱内粗糙的底板上放着一个小木块。木箱和小木块都具有一定的质量。现使木箱获得一个向右的初速度 $v_{0}$，则\xzanswer{B}。（填选项前的字母）

\onepicture[align=right, w=4.5cm]{figs/29e.png}

\fourchoices[answer=B]
{小木块和木箱最终都将静止}
{小木块最终将相对木箱静止，二者一起向右运动}
{小木块在木箱内壁将始终来回往复碰撞，而木箱一直向右运动}
{如果小木块与木箱的左壁碰撞后相对木箱静止，则二者将一起向左运动}
%% answer: B
%% memo:
\memoanswer{
（1）由公式 $m=m_{0}\left(\frac{1}{2}\right)^{\frac{t}{\tau}}$ 可知 C 答案正确。

（2）系统不受外力，系统动量守恒，最终两个物体以相同的速度一起向右运动，B 正确。
}

\end{enumerate}

\end{document}
